\section{Composition Series and Jordan--H\"older Theory}

\begin{definition}
    A finite chain of submodules
    \[
        U=U_0\supsetneq U_1\supsetneq \cdots \supsetneq U_n=0
    \]
    is called a \emph{composition series} of $U$ if each factor
    \[
        U_i/U_{i+1}
    \]
    is simple. The integer $n$ is called the \emph{composition length} of $U$.
\end{definition}

\begin{lemma}
    A nonzero $R$-module $U$ has a composition series if and only if $U$ is both
    Noetherian and Artinian.
\end{lemma}

\begin{proof}
    Suppose first that
    \[
        U=U_0\supsetneq U_1\supsetneq \cdots \supsetneq U_n=0
    \]
    is a composition series. We argue by induction on $n$. If $n=1$, then $U$ is simple,
    hence both Noetherian and Artinian. If $n>1$, then $U_1$ has a composition series of
    length $n-1$, so by induction $U_1$ is Noetherian and Artinian. The quotient $U/U_1$ is simple, hence also Noetherian and Artinian. By the permanence of both
    properties in short exact sequences, $U$ is Noetherian and Artinian.

    Conversely, assume that $U$ is Noetherian and Artinian. Since $U\neq 0$, choose a
    maximal proper submodule $U_1$ of $U$. Then $U/U_1$ is simple. Because $U_1$ is a
    submodule of both a Noetherian and an Artinian module, it is again Noetherian and
    Artinian. Repeating the argument inductively and using the Artinian condition to rule
    out an infinite descending chain, we obtain a composition series.
\end{proof}

\begin{lemma}
    Suppose that
    \[
        U=S_1+\cdots+S_n
    \]
    where each $S_i$ is a simple submodule of $U$. If $V\leq U$, then there exists a
    subset $I\subseteq \{1,\dots,n\}$ such that
    \[
        U=V\oplus \bigoplus_{i\in I} S_i.
    \]
    In particular, every submodule of a finite direct sum of simple modules is a direct
    summand.
\end{lemma}

\begin{proof}
    Choose a subset $I\subseteq \{1,\dots,n\}$ maximal with the property that
    \[
        V\oplus \bigoplus_{i\in I} S_i
    \]
    is a direct sum, and set
    \[
        W:=V\oplus \bigoplus_{i\in I} S_i.
    \]
    We claim that $W=U$. If not, then there exists $j\notin I$ such that
    \[
        S_j\nsubseteq W.
    \]
    Now $W\cap S_j$ is a submodule of the simple module $S_j$, so either
    \[
        W\cap S_j=0
        \qquad \text{or} \qquad
        W\cap S_j=S_j.
    \]
    The second case is impossible because it would imply $S_j\subseteq W$. Hence $W\cap S_j=0$, and therefore
    \[
        W\oplus S_j
    \]
    is a direct sum, contradicting the maximality of $I$. Thus $W=U$.
\end{proof}

\begin{theorem}
    Let $U$ be an $R$-module with a composition series. The following are equivalent.
    \begin{enumerate}
        \item $U$ is semisimple.
        \item $U$ is a sum of simple submodules.
        \item Every submodule of $U$ is a direct summand of $U$.
    \end{enumerate}
\end{theorem}

\begin{proof}
    \textbf{(1) $\Longrightarrow$(2).}
    This is immediate from the definition of a semisimple module.

    \textbf{(2) $\Longrightarrow$(3).}
    Since $U$ has finite composition length, a finite number of simple summands already
    generate $U$. Thus
    \[
        U=S_1+\cdots+S_n
    \]
    for simple submodules $S_i$. The previous lemma shows that every submodule of $U$ is
    a direct summand.

    \textbf{(3) $\Longrightarrow$(1).}
    Let
    \[
        U=U_0\supsetneq U_1\supsetneq \cdots \supsetneq U_n=0
    \]
    be a composition series. Since $U_1$ is a direct summand of $U$, there exists a
    submodule $S_1$ such that
    \[
        U=U_1\oplus S_1.
    \]
    Because $U/U_1$ is simple, the submodule $S_1$ is simple. Repeating the argument
    inside $U_1$, we obtain simple submodules $S_2,\dots,S_n$ with
    \[
        U_i=U_{i+1}\oplus S_{i+1}
    \]
    for all $i$. Hence
    \[
        U=S_1\oplus \cdots \oplus S_n,
    \]
    so $U$ is semisimple.
\end{proof}

\begin{corollary}
    Let $U$ be an $R$-module of finite composition length.
    \begin{enumerate}
        \item The sum of all simple submodules of $U$ is semisimple, hence is the unique
              largest semisimple submodule of $U$.
        \item If $S$ is a simple submodule of $U$, then the sum of all submodules of $U$ isomorphic to $S$ is a submodule of $U$ isomorphic to a direct sum of
              copies of $S$. It is the unique largest submodule of $U$ with this
              property.
    \end{enumerate}
\end{corollary}

\begin{proof}
    Since \(U\) has finite composition length, \(U\) is both Noetherian and Artinian.
    In particular, every submodule of \(U\) is finitely generated.

    \medskip

    \noindent
    \textbf{(1).} Let
    \[
        \Sigma := \sum_{\substack{T\leq U\\ T\text{ simple}}} T
    \]
    be the sum of all simple submodules of \(U\).

    Since \(\Sigma\leq U\) and \(U\) is Noetherian, the submodule \(\Sigma\) is finitely
    generated. Thus there exist elements
    \[
        x_1,\dots,x_m\in \Sigma
    \]
    such that
    \[
        \Sigma=Rx_1+\cdots+Rx_m.
    \]
    By the definition of a sum of submodules, each \(x_j\in\Sigma\) lies in a finite
    sum of simple submodules of \(U\). Hence for each \(j\), there exist simple
    submodules
    \[
        T_{j1},\dots,T_{jr_j}\leq U
    \]
    such that
    \[
        x_j\in T_{j1}+\cdots+T_{jr_j}.
    \]
    Collecting all these finitely many simple submodules, we obtain simple submodules
    \[
        T_1,\dots,T_n\leq U
    \]
    such that
    \[
        x_1,\dots,x_m\in T_1+\cdots+T_n.
    \]
    Therefore
    \[
        \Sigma
        =
        Rx_1+\cdots+Rx_m
        \subseteq
        T_1+\cdots+T_n
        \subseteq
        \Sigma.
    \]
    Hence
    \[
        \Sigma=T_1+\cdots+T_n.
    \]

    Since \(\Sigma\) is a finite sum of simple submodules, by the previous theorem
    \(\Sigma\) is semisimple.

    Now let \(W\leq U\) be any semisimple submodule. Then \(W\) is a sum of simple
    submodules of \(W\). But every simple submodule of \(W\) is also a simple submodule
    of \(U\). Hence every simple summand appearing in \(W\) is contained in the defining
    sum of \(\Sigma\). Therefore
    \[
        W\subseteq \Sigma.
    \]
    Thus \(\Sigma\) contains every semisimple submodule of \(U\), and so \(\Sigma\) is the
    unique largest semisimple submodule of \(U\).

    \medskip

    \noindent
    \textbf{(2).} Let \(S\leq U\) be a simple submodule, and define
    \[
        \Sigma_S
        :=
        \sum_{\substack{T\leq U\\ T\cong S}} T.
    \]
    This is the sum of all submodules of \(U\) which are isomorphic to \(S\).

    Since \(\Sigma_S\leq U\) and \(U\) is Noetherian, \(\Sigma_S\) is finitely generated.
    Thus there exist elements
    \[
        y_1,\dots,y_\ell\in \Sigma_S
    \]
    such that
    \[
        \Sigma_S=Ry_1+\cdots+Ry_\ell.
    \]
    Again, by the definition of a sum of submodules, each \(y_j\) lies in a finite sum
    of submodules of \(U\) isomorphic to \(S\). Hence, after collecting all the finitely
    many such submodules involved, there exist submodules
    \[
        S_1,\dots,S_m\leq U
    \]
    with
    \[
        S_i\cong S
        \qquad
        \text{for all } i,
    \]
    such that
    \[
        \Sigma_S=S_1+\cdots+S_m.
    \]

    Each \(S_i\) is simple, since it is isomorphic to the simple module \(S\). Hence
    \(\Sigma_S\) is a finite sum of simple submodules. By the previous theorem,
    \(\Sigma_S\) is semisimple. More precisely, applying the preceding lemma to the
    finite sum
    \[
        \Sigma_S=S_1+\cdots+S_m
    \]
    with \(V=0\), we may choose a subset \(I\subseteq\{1,\dots,m\}\) such that
    \[
        \Sigma_S=\bigoplus_{i\in I} S_i.
    \]
    Since each \(S_i\cong S\), we obtain
    \[
        \Sigma_S\cong S^{\oplus |I|}.
    \]
    Thus \(\Sigma_S\) is isomorphic to a direct sum of copies of \(S\).

    Finally, let \(W\leq U\) be any submodule which is isomorphic to a direct sum of
    copies of \(S\). Then \(W\) is the sum of its simple submodules, each of which is
    isomorphic to \(S\). Since these simple submodules are also submodules of \(U\),
    they are all contained in the defining sum of \(\Sigma_S\). Hence
    \[
        W\subseteq \Sigma_S.
    \]
    Therefore \(\Sigma_S\) contains every submodule of \(U\) which is isomorphic to a
    direct sum of copies of \(S\). Hence \(\Sigma_S\) is the unique largest submodule of
    \(U\) with this property.
\end{proof}

\begin{theorem}[Zassenhaus lemma]
    Let $A'\subseteq A$ and $B'\subseteq B$ be submodules of an $R$-module $M$.
    Then
    \[
        \frac{A'+(A\cap B)}{A'+(A\cap B')}
        \cong
        \frac{B'+(A\cap B)}{B'+(A'\cap B)}.
    \]
\end{theorem}

\begin{proof}
    Set
    \[
        C:=(A\cap B')+(A'\cap B).
    \]
    We shall show that both quotient modules are naturally isomorphic to
    \[
        \frac{A\cap B}{C}.
    \]

    First consider the quotient
    \[
        \frac{A'+(A\cap B)}{A'+(A\cap B')}.
    \]
    Apply the second isomorphism theorem to the submodules
    \[
        X:=A\cap B,
        \qquad
        Y:=A'+(A\cap B')
    \]
    of $M$. Since
    \[
        X+Y=(A\cap B)+A'+(A\cap B')=A'+(A\cap B),
    \]
    we obtain
    \[
        \frac{A'+(A\cap B)}{A'+(A\cap B')}
        =
        \frac{X+Y}{Y}
        \cong
        \frac{X}{X\cap Y}.
    \]
    It remains to compute $X\cap Y$. We claim that
    \[
        (A\cap B)\cap \bigl(A'+(A\cap B')\bigr)
        =
        (A'\cap B)+(A\cap B').
    \]
    The inclusion ``$\supseteq$'' is immediate, because
    \[
        A'\cap B\subseteq A\cap B,
        \qquad
        A\cap B'\subseteq A\cap B,
    \]
    and both summands are contained in $A'+(A\cap B')$.

    Conversely, let
    \[
        x\in (A\cap B)\cap \bigl(A'+(A\cap B')\bigr).
    \]
    Then $x\in A\cap B$, and we may write
    \[
        x=a'+c
    \]
    with $a'\in A'$ and $c\in A\cap B'$. Since $x\in B$ and $c\in B'\subseteq B$, we have
    \[
        a'=x-c\in B.
    \]
    Hence
    \[
        a'\in A'\cap B.
    \]
    Therefore
    \[
        x=a'+c\in (A'\cap B)+(A\cap B').
    \]
    This proves the claimed equality. Consequently,
    \[
        \frac{A'+(A\cap B)}{A'+(A\cap B')}
        \cong
        \frac{A\cap B}{(A'\cap B)+(A\cap B')}.
    \]

    Similarly, consider the quotient
    \[
        \frac{B'+(A\cap B)}{B'+(A'\cap B)}.
    \]
    Apply the second isomorphism theorem to
    \[
        X:=A\cap B,
        \qquad
        Y:=B'+(A'\cap B).
    \]
    Then
    \[
        X+Y=(A\cap B)+B'+(A'\cap B)=B'+(A\cap B),
    \]
    so
    \[
        \frac{B'+(A\cap B)}{B'+(A'\cap B)}
        \cong
        \frac{A\cap B}{(A\cap B)\cap \bigl(B'+(A'\cap B)\bigr)}.
    \]
    We now compute the intersection. We claim that
    \[
        (A\cap B)\cap \bigl(B'+(A'\cap B)\bigr)
        =
        (A\cap B')+(A'\cap B).
    \]
    Again the inclusion ``$\supseteq$'' is clear. Conversely, let
    \[
        x\in (A\cap B)\cap \bigl(B'+(A'\cap B)\bigr).
    \]
    Then $x\in A\cap B$, and we may write
    \[
        x=b'+d
    \]
    with $b'\in B'$ and $d\in A'\cap B$. Since $x\in A$ and $d\in A'\subseteq A$, we get
    \[
        b'=x-d\in A.
    \]
    Hence
    \[
        b'\in A\cap B'.
    \]
    Therefore
    \[
        x=b'+d\in (A\cap B')+(A'\cap B).
    \]
    This proves the second intersection equality. Thus
    \[
        \frac{B'+(A\cap B)}{B'+(A'\cap B)}
        \cong
        \frac{A\cap B}{(A\cap B')+(A'\cap B)}.
    \]

    We have shown that both quotient modules are naturally isomorphic to the same
    quotient
    \[
        \frac{A\cap B}{(A\cap B')+(A'\cap B)}.
    \]
    Hence
    \[
        \frac{A'+(A\cap B)}{A'+(A\cap B')}
        \cong
        \frac{B'+(A\cap B)}{B'+(A'\cap B)}.
    \]
\end{proof}

\begin{definition}
    Let
    \[
        U=U_0\supseteq U_1\supseteq\cdots\supseteq U_m=0
    \]
    be a finite submodule series of an $R$-module $U$.
    A refinement of this series is another finite submodule series obtained by inserting
    additional submodules between the terms $U_{i-1}$ and $U_i$.

    Two finite submodule series are said to be equivalent if their nonzero successive
    quotient modules are isomorphic up to a permutation.
\end{definition}

\begin{theorem}[Schreier refinement theorem]
    Any two finite submodule series of an $R$-module admit equivalent refinements.
\end{theorem}

\begin{proof}
    Let the two finite submodule series be
    \[
        U=U_0\supseteq U_1\supseteq\cdots\supseteq U_m=0
    \]
    and
    \[
        U=V_0\supseteq V_1\supseteq\cdots\supseteq V_n=0.
    \]
    We shall construct refinements of both series and then compare their successive
    quotient modules.

    For each $1\leq i\leq m$ and $0\leq j\leq n$, define
    \[
        U_{i,j}:=U_i+(U_{i-1}\cap V_j).
    \]
    Since
    \[
        V_0=U
        \qquad\text{and}\qquad
        V_n=0,
    \]
    we have
    \[
        U_{i,0}
        =
        U_i+(U_{i-1}\cap V_0)
        =
        U_i+U_{i-1}
        =
        U_{i-1},
    \]
    and
    \[
        U_{i,n}
        =
        U_i+(U_{i-1}\cap V_n)
        =
        U_i.
    \]
    Moreover, since
    \[
        V_0\supseteq V_1\supseteq\cdots\supseteq V_n,
    \]
    we get
    \[
        U_{i,0}\supseteq U_{i,1}\supseteq\cdots\supseteq U_{i,n}.
    \]
    Thus, for each fixed $i$, the chain
    \[
        U_{i-1}=U_{i,0}\supseteq U_{i,1}\supseteq\cdots\supseteq U_{i,n}=U_i
    \]
    refines the interval
    \[
        U_{i-1}\supseteq U_i.
    \]
    Splicing these chains together, we obtain a refinement of the first series:
    \[
        U=U_{1,0}\supseteq U_{1,1}\supseteq\cdots\supseteq U_{1,n}
        =
        U_1=U_{2,0}\supseteq\cdots\supseteq U_{m,n}=0.
    \]

    Similarly, for each $1\leq j\leq n$ and $0\leq i\leq m$, define
    \[
        V_{j,i}:=V_j+(V_{j-1}\cap U_i).
    \]
    Then
    \[
        V_{j,0}
        =
        V_j+(V_{j-1}\cap U_0)
        =
        V_j+V_{j-1}
        =
        V_{j-1},
    \]
    and
    \[
        V_{j,m}
        =
        V_j+(V_{j-1}\cap U_m)
        =
        V_j.
    \]
    Hence, for each fixed $j$, we have a chain
    \[
        V_{j-1}=V_{j,0}\supseteq V_{j,1}\supseteq\cdots\supseteq V_{j,m}=V_j.
    \]
    Splicing these chains together gives a refinement of the second series:
    \[
        U=V_{1,0}\supseteq V_{1,1}\supseteq\cdots\supseteq V_{1,m}
        =
        V_1=V_{2,0}\supseteq\cdots\supseteq V_{n,m}=0.
    \]

    We now compare the successive quotients of these two refinements.
    Fix indices
    \[
        1\leq i\leq m,
        \qquad
        1\leq j\leq n.
    \]
    Apply the Zassenhaus lemma to the four submodules
    \[
        A'=U_i,\qquad A=U_{i-1},
        \qquad
        B'=V_j,\qquad B=V_{j-1}.
    \]
    Since $U_i\subseteq U_{i-1}$ and $V_j\subseteq V_{j-1}$, the Zassenhaus lemma gives
    \[
        \frac{U_i+(U_{i-1}\cap V_{j-1})}
             {U_i+(U_{i-1}\cap V_j)}
        \cong
        \frac{V_j+(U_{i-1}\cap V_{j-1})}
             {V_j+(U_i\cap V_{j-1})}.
    \]
    In terms of the notation introduced above, this is precisely
    \[
        \frac{U_{i,j-1}}{U_{i,j}}
        \cong
        \frac{V_{j,i-1}}{V_{j,i}}.
    \]

    Therefore every successive quotient in the refinement of the first series is
    isomorphic to a corresponding successive quotient in the refinement of the second
    series. The correspondence is given by the pairs of indices
    \[
        (i,j)\longleftrightarrow (j,i).
    \]
    Hence the two constructed refinements are equivalent.

    If some consecutive terms in the constructed refinements coincide, then the
    corresponding quotient is zero. Deleting these repeated terms on both sides does not
    affect the list of nonzero quotient modules. Thus the two original finite submodule
    series admit equivalent refinements.
\end{proof}

\begin{theorem}[Jordan--H\"older theorem]
    Let $U$ be an $R$-module with a composition series. Then any two composition
    series of $U$ have the same length, and their composition factors are the same up to
    permutation and isomorphism.
\end{theorem}

\begin{proof}
    Let
    \[
        U=U_0\supsetneq U_1\supsetneq\cdots\supsetneq U_m=0
    \]
    and
    \[
        U=V_0\supsetneq V_1\supsetneq\cdots\supsetneq V_n=0
    \]
    be two composition series of $U$. Thus each quotient
    \[
        U_{i-1}/U_i
    \]
    and each quotient
    \[
        V_{j-1}/V_j
    \]
    is simple.

    By the Schreier refinement theorem, these two finite submodule series admit
    equivalent refinements. We now observe that a composition series has no proper
    refinement.

    Indeed, suppose that one could insert a submodule $W$ strictly between two
    consecutive terms of the first composition series:
    \[
        U_{i-1}\supsetneq W\supsetneq U_i.
    \]
    Then $W/U_i$ would be a nonzero proper submodule of
    \[
        U_{i-1}/U_i.
    \]
    This contradicts the fact that $U_{i-1}/U_i$ is simple. Hence no such intermediate
    submodule $W$ exists. Therefore any refinement of the first composition series only
    inserts repeated terms. After deleting repetitions, the refinement is just the
    original composition series. The same argument applies to the second composition
    series.

    Since the Schreier refinements of the two composition series are equivalent, and
    since those refinements reduce back to the original composition series after deleting
    repeated terms, the original composition series themselves are equivalent.

    Therefore the number of nonzero successive quotient modules in the two series is
    the same. Hence
    \[
        m=n.
    \]
    Moreover, there exists a permutation $\sigma\in S_m$ such that
    \[
        U_{i-1}/U_i
        \cong
        V_{\sigma(i)-1}/V_{\sigma(i)}
    \]
    for every $1\leq i\leq m$.

    Thus any two composition series of $U$ have the same length and the same
    composition factors up to permutation and isomorphism.
\end{proof}

